% =====================================================================
\chapter{Autovalores, autovectores y teorema espectral}
\label{tema:3}
% =====================================================================

Un autovector de una matriz cuadrada es una dirección que la matriz no gira, solo estira o encoge. Este tema define autovalores y autovectores, estudia la diagonalización y demuestra el resultado central para las aplicaciones: las matrices reales simétricas tienen autovalores reales y una base ortonormal de autovectores. Se sigue a Strang~\cite[§6.1--6.3]{strang-ila}, las clases~21, 22 y~25 de 18.06~\cite{mit1806}, las clases~4 y~5 de 18.065~\cite{mit18065} y los §4.2 y~§4.4 de~\cite{mml}.

\section{Definición y polinomio característico}

\begin{definicion}[Autovalor y autovector]
Sea $A\in\R^{n\times n}$. Un número $\lambda\in\mathbb{C}$ es un \emph{autovalor} de $A$ si existe $v\ne 0$ tal que
\[
Av=\lambda v.
\]
Todo $v\ne 0$ que cumple la igualdad es un \emph{autovector} asociado a $\lambda$. El conjunto $\Nuc(A-\lambda I)$, formado por los autovectores y el vector nulo, es el \emph{espacio propio} de $\lambda$.
\end{definicion}

\begin{proposicion}[Polinomio característico]\label{prop:polcar}
$\lambda$ es autovalor de $A$ si y solo si $\det(A-\lambda I)=0$. La función $p_A(\lambda)=\det(A-\lambda I)$ es un polinomio de grado $n$ en $\lambda$, llamado \emph{polinomio característico}.
\end{proposicion}

Por el teorema fundamental del álgebra, $p_A$ tiene $n$ raíces complejas contadas con multiplicidad: $p_A(\lambda)=\prod_{i=1}^{n}(\lambda_i-\lambda)$.

\addtocounter{definicion}{1}%

\begin{proposicion}[Traza y determinante]\label{prop:traza-det}
Si $\lambda_1,\dots,\lambda_n$ son los autovalores de $A$ contados con multiplicidad, entonces
\[
\det(A)=\prod_{i=1}^{n}\lambda_i,\qquad \traza(A)=\sum_{i=1}^{n}a_{ii}=\sum_{i=1}^{n}\lambda_i.
\]
\end{proposicion}

En el \cref{ej:S}, $3+1=4=\traza(S)$ y $3\cdot 1=3=\det(S)$. Estas dos comprobaciones detectan la mayoría de los errores de cálculo.

\begin{table}[h]
\centering\small
\caption{Autovalores de algunas matrices de orden 2.}\label{tab:autovalores}
\begin{tabular}{@{}lcl@{}}
\toprule
Matriz & Autovalores & Interpretación \\
\midrule
$\begin{pmatrix}2&1\\1&2\end{pmatrix}$ & $3$ y $1$ & estira por~3 la diagonal y conserva la antidiagonal \\[6pt]
$\begin{pmatrix}3&1\\0&2\end{pmatrix}$ & $3$ y $2$ & triangular: los autovalores son la diagonal \\[6pt]
$\begin{pmatrix}1&0\\0&0\end{pmatrix}$ & $1$ y $0$ & proyección: conserva el eje $x$ y anula el eje $y$ \\[6pt]
$\begin{pmatrix}0&-1\\1&0\end{pmatrix}$ & $\pm i$ & giro de $90^{\circ}$: ninguna dirección real se conserva \\
\bottomrule
\end{tabular}
\end{table}

\section{Diagonalización}

\begin{proposicion}\label{prop:indep-autovectores}
Autovectores asociados a autovalores distintos son linealmente independientes.
\end{proposicion}

\begin{teorema}[Diagonalización]\label{teo:diagonalizacion}
$A\in\R^{n\times n}$ admite una factorización $A=X\Lambda X^{-1}$, con $\Lambda$ diagonal y $X$ invertible, si y solo si $A$ tiene $n$ autovectores linealmente independientes. En ese caso, las columnas de $X$ son autovectores, la diagonal de $\Lambda$ contiene los autovalores correspondientes y $A^{k}=X\Lambda^{k}X^{-1}$ para todo $k\ge 1$.
\end{teorema}

Por la \cref{prop:indep-autovectores}, una matriz con $n$ autovalores distintos es diagonalizable. Con autovalores repetidos puede no serlo: $\begin{pmatrix}1&1\\0&1\end{pmatrix}$ tiene el autovalor doble~1, pero su espacio propio es $\gen\{e_1\}$, de dimensión~1. La diagonalizabilidad no guarda relación con la invertibilidad (\cref{eje:3-mml}).

\section{Matrices simétricas y teorema espectral}

\begin{teorema}\label{teo:simetricas}
Sea $S\in\R^{n\times n}$ simétrica, $S\T=S$. Entonces:
\begin{enumerate}
  \item todos los autovalores de $S$ son reales;
  \item autovectores asociados a autovalores distintos son ortogonales.
\end{enumerate}
\end{teorema}

\begin{teorema}[Teorema espectral]\label{teo:espectral}
Toda matriz real simétrica $S\in\R^{n\times n}$ se factoriza como
\[
S=Q\Lambda Q\T=\sum_{i=1}^{n}\lambda_i\,q_iq_i\T,
\]
donde $Q=[q_1\mid\cdots\mid q_n]$ es ortogonal y $\Lambda=\diag(\lambda_1,\dots,\lambda_n)$ es real.
\end{teorema}

Leída de derecha a izquierda, $S=Q\Lambda Q\T$ dice que $S$ expresa el vector en la base de autovectores ($Q\T$), estira cada eje por su autovalor ($\Lambda$) y vuelve a la base original ($Q$). Para el \cref{ej:S},
\[
S=3\,q_1q_1\T+1\,q_2q_2\T=\frac{3}{2}\begin{pmatrix}1&1\\1&1\end{pmatrix}+\frac{1}{2}\begin{pmatrix}1&-1\\-1&1\end{pmatrix}.
\]

\begin{definicion}[Matriz semidefinida positiva]
Una matriz simétrica $S$ es \emph{semidefinida positiva} si $x\T Sx\ge 0$ para todo $x$, y \emph{definida positiva} si $x\T Sx>0$ para todo $x\ne 0$.
\end{definicion}

\begin{proposicion}\label{prop:AtA}
Para toda $A\in\R^{m\times n}$, la matriz $A\T A$ es simétrica y semidefinida positiva, y sus autovalores son mayores o iguales que cero.
\end{proposicion}

\begin{proposicion}[Cociente de Rayleigh]\label{prop:rayleigh}
Sea $S$ simétrica con autovalores $\lambda_1\ge\dots\ge\lambda_n$ y autovectores ortonormales $q_1,\dots,q_n$. Entonces
\[
\max_{\norm{x}=1}x\T Sx=\lambda_1,
\]
y el máximo se alcanza en $x=q_1$.
\end{proposicion}

Este resultado es la base del análisis de componentes principales (\cref{tema:5}): la dirección de máxima varianza de unos datos es el primer autovector de su matriz de covarianza.

\section{El método de la potencia}

Si $S$ es simétrica con $\abs{\lambda_1}>\abs{\lambda_2}\ge\dots\ge\abs{\lambda_n}$, la iteración
\[
x_{k+1}=\frac{Sx_k}{\norm{Sx_k}}
\]
converge, salvo signo, al autovector $q_1$ siempre que $x_0$ tenga componente no nula sobre $q_1$. El error decrece como $\abs{\lambda_2/\lambda_1}^{k}$ (\cref{eje:3-potencia}).

\addtocounter{definicion}{1}%

\addtocounter{definicion}{1}%

